\clearpage\SourceLeaf{159}%
\thispagestyle{empty}
\pdfbookmark[0]{Anzeigen eigner Schriften}{notices-title}
\begin{center}
\vspace*{1in}{\Large ANZEIGEN\par}
\vspace{.4in}{\large EIGNER\par}
\vspace{.4in}{\Large SCHRIFTEN.\par}
\end{center}

\clearpage\SourceLeaf{161}%
\renewcommand{\BookRunningTitle}{Anzeigen eigner Schriften}
\pdfbookmark[1]{Theorematis arithmetici demonstratio nova. Anzeige 1808 Mai 12}{notice-1808-05-12}
\begin{center}
\rule{.86\linewidth}{.3pt}\par
Göttingische gelehrte Anzeigen.\quad 1808 Mai 12.\par
\rule{.86\linewidth}{.3pt}
\end{center}

Eine vom Herrn Prof. \textsc{Gauss} am 15. Januar d. J. der königl. Societät der Wissenschaften überreichte Abhandlung,
\begin{center}\emph{Theorematis arithmetici demonstratio nova},\end{center}
deren Inhaltsanzeige wir hier noch nachzuholen haben, hat das berühmte Fundamental-Theorem der Lehre von den quadratischen Resten zum Gegenstande, welches sowohl in der ganzen \emph{höhern Arithmetik}, als in den angrenzenden Theilen der Analysis eine so wichtige Rolle spielt. Bekanntlich heisst eine ganze Zahl $a$ \emph{quadratischer Rest} der ganzen Zahl $b$, wenn es Zahlen der Form $xx-a$ gibt, die durch $b$ theilbar sind, sowie im entgegengesetzten Falle $a$ \emph{quadratischer Nichtrest} von $b$ genannt wird: die Zahl $a$ kann positiv oder negativ sein, $b$ hingegen wird immer als positiv angesehen. Die höhere Arithmetik lehrt, dass alle Primzahlen $b$, für welche eine gegebene Zahl $a$ quadratischer Rest ist, unter gewissen linearischen Formen begriffen sind, so wie wiederum andere linearische Formen alle Primzahlen enthalten, von denen $a$ Nichtrest ist. So ist z. B. $-1$ quadratischer Rest aller Primzahlen der Form $4n+1$, quadratischer Nichtrest aller Primzahlen der Form $4n+3$; ferner $+2$ ist quadratischer Rest aller Primzahlen der Form $8n+1$, $8n+7$, hingegen quadratischer Nichtrest aller Primzahlen der Formen $8n+3$, $8n+5$. Aehnlicher specieller Lehrsätze gibt es eine unendliche Menge, die sich aber alle aus der Verbindung der beiden angeführten
\SourceLeaf{162}%
mit folgendem allgemeinen ableiten lassen: Zwei ungleiche positive (ungerade) Primzahlen, $p$, $q$, haben allemal \emph{gleiche} Relation wechselseitig zu einander (d. i. die eine ist quadratischer Rest oder Nichtrest der andern, je nachdem die andere Rest oder Nichtrest der ersten ist), wenn entweder beide von der Form $4n+1$ sind, oder wenigstens die eine: hingegen ist ihre wechselseitige Relation entgegengesetzt (d. i. die eine ist Nichtrest der andern, wenn diese Rest von jener ist, und umgekehrt), so oft \emph{beide} zugleich von der Form $4n+3$ sind. Dies ist das erwähnte Fundamental-Theorem, welches man in mehr als einer Gestalt ausdrücken kann: die hier gewählte ist diejenige, in der es in der Abhandlung des Hrn. Prof. \textsc{Gauss} neu bewiesen ist.

Die schönsten Lehrsätze der höhern Arithmetik, und namentlich auch diejenigen, wovon hier die Rede ist, haben das Eigne, dass sie durch Induction leicht entdeckt werden, ihre Beweise hingegen äusserst versteckt liegen, und nur durch sehr tief eindringende Untersuchungen aufgespürt werden können. Gerade dies ist es, was der höhern Arithmetik jenen zauberischen Reiz gibt, der sie zur Lieblingswissenschaft der ersten Geometer gemacht hat, ihres unerschöpflichen Reichthums nicht zu gedenken, woran sie alle andere Theile der reinen Mathematik so weit übertrifft. Die beiden oben erwähnten Specialsätze waren schon \textsc{Fermat} bekannt, welcher, seiner Behauptung nach, auch im Besitz ihrer Beweise war: ob er sich darin nicht täuschte, können wir nicht entscheiden, da er nie Etwas davon bekannt gemacht hat: aber für möglich dürfen wir es gewiss halten, da mehrere Beispiele von Selbsttäuschung bei andern grossen Geometern, namentlich bei \textsc{Euler}, \textsc{Legendre} und auch bei \textsc{Fermat} selbst, vorhanden sind. Von dem ersten jener Theoreme gab \textsc{Euler} den ersten Beweis; allein das andere zu demonstriren, glückte diesem grossen Geometer, seiner eifrigen, viele Jahre hindurch fortgesetzten, Bemühungen ungeachtet, nicht; erst \textsc{Lagrange} war es vorbehalten, diese Lücke auszufüllen. Beide Geometer bewiesen auch noch verschiedene andere specielle Sätze, eine grössere Anzahl aber, die sie durch Induction fanden, entzog sich ihren Bemühungen, sie zu beweisen, stets. Es ist indess ein merkwürdiges Spiel des Zufalls, dass beide Geometer durch Induction nicht auf das allgemeine Fundamental-Theorem gekommen sind, das einer so einfachen Darstellung fähig ist. Dieses ist zuerst, obwohl in einer etwas andern Gestalt, von \textsc{Legendre} vorgetragen, in der \emph{Histoire de l'Académie des Sciences de Paris} 1785; sowohl hier, als nachher in seinem Werke: \emph{Essai d'une théorie des nombres}, hat
\SourceLeaf{163}%
dieser treffliche Analyst den Beweis auf sehr scharfsinnige Untersuchungen zu gründen gesucht, die aber gleichwohl nicht zu dem gewünschten Ziele geführt haben, welches, wenn wir uns nicht irren, auch auf diesem Wege nicht erreicht werden konnte.

Der Verfasser der Abhandlung, welcher diese Anzeige gewidmet ist, betrat die Bahn der höhern Arithmetik zu einer Zeit, wo ihm alle frühern Arbeiten andrer Geometer in dieser Wissenschaft ganz unbekannt waren: diesem Umstande ist es hauptsächlich zuzuschreiben, dass er überall einen ganz eigenthümlichen Gang genommen hat. Jenes Fundamental-Theorem fand er zwar schon sehr früh durch Induction, allein erst ein ganzes Jahr später gelang es ihm, nach vielen Schwierigkeiten und vergeblichen Versuchen, den ersten vollkommen strengen Beweis aufzufinden, der im vierten Abschnitte seiner \emph{Disquisitiones Arithmeticae} entwickelt ist: dieser Beweis gründet sich aber auf sehr mühsame und weitläufige Auseinandersetzungen. In der Folge kam er noch auf drei andre Beweise, die zwar von jener Unbequemlichkeit frei sind, aber dagegen andre sehr tiefliegende und ihrem Inhalte nach ganz heterogene Untersuchungen voraussetzen: der eine dieser Beweise ist gleichfalls in dem angeführten Werke Art. 262 mitgetheilt, die beiden andern werden zu ihrer Zeit bekannt gemacht werden. Immer blieb also noch der Wunsch übrig, dass es möglich sein möchte, einen kürzern, von fremdartigen Untersuchungen unabhängigen, Beweis zu entdecken. Der Verf. hofft daher, dass die Freunde der höhern Arithmetik mit Vergnügen einen fünften Beweis sehen werden, der in gegenwärtiger Abhandlung auf weniger als fünf Seiten vorgetragen ist, und in jeder Hinsicht nichts zu wünschen übrig zu lassen scheint. Bei der gedrängten Kürze, worin dieser Beweis abgefasst ist, können wir freilich hier von dem Gange desselben nur eine unvollkommene Idee geben: mehr würde hier aber auch um so überflüssiger sein, da der XVIte Band der \emph{Commentationes}, worin er bereits abgedruckt ist, nächstens erscheinen wird.

Die Grundlage des Beweises ist folgender neuer Lehrsatz: Wenn $p$ eine (positive ungerade) Primzahl, $k$ eine beliebige, durch $p$ nicht theilbare, ganze Zahl bedeutet; wenn ferner unter den Resten, die aus der Division der $\frac12(p-1)$ Producte $k$, $2k$, $3k\dotsfour\frac12(p-1)k$ durch $p$ entstehen, in allen sich $\mu$ Reste befinden, die grösser als $\frac12p$ sind (also $\frac12(p-1)-\mu$ solche, die kleiner sind als $\frac12p$), so wird $k$ ein quadratischer Rest von $p$ sein, wenn $\mu$ gerade ist, hingegen ein quadratischer Nichtrest, wenn $\mu$ ungerade ist. Die Zahl $\mu$, die bloss von $k$
\SourceLeaf{164}%
und $p$ abhängig ist, mag durch das Zeichen $(k,p)$ dargestellt werden. Durch eine Reihe von Schlüssen, die keines Auszugs fähig sind, wird nun gezeigt, dass, wenn $k$ und $p$ zwei ungerade Zahlen sind, die keinen gemeinschaftlichen Theiler haben, allemal $(k,p)+(p,k)+\frac14(k-1)(p-1)$ eine \emph{gerade} Zahl wird: daraus folgt also, dass, so oft $k$ und $p$ beide von der Form $4n+3$ sind, nothwendig eine der Zahlen $(k,p)$, $(p,k)$ gerade, die andere ungerade sein muss; in allen übrigen Fällen hingegen, d. i. so oft beiden Zahlen $k$ und $p$, oder wenigstens einer, die Form $4n+1$ zukommt, werden nothwendig entweder .$(k,p)$, $(p,k)$ beide zugleich gerade, oder beide zugleich ungerade sein. Hieraus folgt, in Verbindung mit obigem Lehrsatze, die Wahrheit des Fundamental-Theorems von selbst.--- Auf demselben Wege, auf dem diese Resultate gefunden werden, wird in der Abhandlung zugleich ein neuer Beweis für die oben erwähnten beiden Specialsätze gegeben: es lässt sich nemlich leicht zeigen, dass $(-1,p)=\frac12(p-1)$, also gerade oder ungerade, je nachdem $p$ die Form $4n+1$ oder $4n+3$ hat; eben so wird $(2,p)=\frac14(p-1)$, wenn $p$ die Form $4n+1$ hat, und $(2,p)=\frac14(p+1)$, wenn $p$ von der Form $4n+3$ ist, daher $(2,p)$ gerade wird, so oft $p$ die Form $8n+1$ oder $8n+7$ hat, hingegen ungerade, so oft $p$ von der Form $8n+3$ oder $8n+5$ ist.
\BookRule{.28\linewidth}

